\chapter{FDTD 与 FEM 在 PCSEL 建模中的作用}
\label{ch:fdtd-fem}

\section*{学习目标}
\begin{itemize}
  \item 理解为什么在 RCWA 之后仍然需要 FDTD 与 FEM，以及它们分别擅长解决哪类 PCSEL 问题。
  \item 掌握几何设置、边界条件、激励源、结果提取和收敛检查的基本逻辑，而不是停留在软件按钮层面。
  \item 清楚地区分冷腔分析、带损耗分析和加入增益后的分析在数学意义上有什么不同。
\end{itemize}

\section*{先修知识}
建议已掌握 \cref{ch:rcwa,ch:threshold,ch:polarization} 中关于开放边界、损耗、远场和阈值的基础概念。

\section*{本章路线图}
上一章的 RCWA 解决的是周期单胞开放结构问题，但真实 PCSEL 最终总要回到有限尺寸器件。边缘、有限孔阵、非理想窗口、局部几何细节和复杂外延层堆栈，都会把问题推向更一般的全波求解。本章先解释为什么 FDTD 和 FEM 会在这里登场，再分别说明它们的控制方程和数值哲学，随后讨论边界条件、激励源、共振参数提取与远场计算，最后严格说明它们能给出什么、又不能单独给出什么。

\section{为什么有限尺寸一旦重要，全波方法就不可回避}
对 PCSEL 来说，真正让单胞方法失去垄断地位的，不只是“器件变大了”，而是有限尺寸带来的物理变了。无限周期模型里，模式由 Bloch 波矢标记；有限器件里，模式必须同时满足边缘终止、横向包络和开放边界条件。于是，边缘泄漏、横向模式竞争、注入窗口截断和热透镜效应都会显式出现。

这时候，如果还只依赖周期单胞方法，就会把有限器件里真正决定模式排序的一部分物理漏掉。FDTD 与 FEM 的意义，正是在于它们不再默认平面内无限周期，而是直接对给定有限几何上的 Maxwell 方程做求解。也正因为如此，它们成为从“单胞光学”走向“真实器件光学”的关键台阶。

\begin{IntuitionBox}
RCWA 擅长告诉你“无限周期单胞在开放边界下会怎样散射”；FDTD 和 FEM 则开始回答“这个真实大小、真实边界、真实窗口的器件会怎样支持模式并向外辐射”。这两层都属于光学，但问题层次已经不同。
\end{IntuitionBox}

\section{同一个 Maxwell 方程，为什么会分化出两种数值哲学}
无论 FDTD 还是 FEM，最终处理的都是 Maxwell 方程。若写成频域形式，可用
\begin{equation}
\nabla\times\mu_r^{-1}\nabla\times\E
-
k_0^2\varepsilon_r\E
=
\ii\omega\mu_0\J,
\label{eq:fem_wave_eq}
\end{equation}
其中 $k_0=\omega/c$，$\J$ 是外加源项。FEM 直接对 \cref{eq:fem_wave_eq} 的弱形式进行离散；FDTD 则回到时域，对 Maxwell 旋度方程逐时步推进。两者的根本差异，不是“谁更高级”，而是看问题的方式不同。

FDTD 通过时间推进观察场如何建立、衰减和传播。它非常适合宽带激励和瞬态过程，因为一次时域仿真可以在后处理时取出多个频率分量。

FEM 则更像是频域定点瞄准。它适合在指定频率附近做高精度场分布分析，也适合直接构造本征值问题来求开放腔的复频率与模场。

简言之，FDTD 更偏“看全过程”，FEM 更偏“精确看某个点或某个模”。

\section{FDTD 的强项：宽带、瞬态与模建立过程}
FDTD 的核心是把时域 Maxwell 方程写在交错的 Yee 网格上，用电场和磁场在半时间步上交替更新。这个方法对 PCSEL 特别有吸引力，主要因为三个原因。

第一，它天然适合宽带问题。用一个短脉冲或调制源激励器件，就可以在一次仿真中得到一段频谱中的响应。这对于寻找候选共振频率和观察多模并存非常方便。

第二，它能够直接看模如何建立与衰减。如果想判断某个共振是快速泄漏还是长寿命，时域包络会给出非常直观的图像。

第三，它容易和时域监视器结合，提取瞬态远场、延迟响应或短时干涉现象。这对某些动态调制或脉冲工作问题尤其有用。

但 FDTD 的代价也很明确。为了保证精度，它需要足够细的空间网格和满足稳定条件的时间步长；三维器件一旦横向尺寸变大，存储和时间成本会迅速上升。PCSEL 又往往同时具有微细孔结构和较大横向尺寸，这正是 FDTD 最容易变得昂贵的原因。

\section{FEM 的强项：复杂几何、本征模与局域高精度}
FEM 的优势在另一边。它通过构造弱形式，把几何复杂性和材料非均匀性更自然地吸收到网格与形函数中。对于 PCSEL 来说，这带来几个实际好处。

第一，复杂孔形、锥形刻蚀、圆角、分段金属和非规则边界都能更自然地表示，而不需要被强行贴合到正交网格。

第二，FEM 非常适合做本征频率求解。通过加入 PML 等开放边界，直接求解复本征频率，就能得到冷腔共振频率和辐射 $Q$。

第三，FEM 在局域场细节分析上往往更强。若你需要精确比较量子阱区、掺杂区、金属邻近区的场重叠，或者精确获取某一支模式的近场矢量结构，FEM 常常更有优势。

当然，FEM 也不是没有代价。它的前处理通常更重，对网格质量和求解器参数更敏感；而且如果目标是扫很宽的频段，频域逐点扫描会比 FDTD 更耗时间。

\begin{RigorousBox}
FDTD 与 FEM 的区别，不是“一个近似，一个精确”。两者都是数值离散方法，都需要收敛检查和边界审计。真正的区别在于：FDTD 的误差更多地表现在时空离散和长时间积分上，FEM 的误差更多地表现在网格质量、弱形式离散和开放边界本征值的谱污染上。
\end{RigorousBox}

\section{几何设置为什么不能只靠“画出来”}
全波方法里，几何设置从来不是一个机械步骤。对 PCSEL 而言，几何模型至少同时承载三类信息：
\begin{itemize}
  \item 晶格和孔形决定的面内反馈与辐射通道；
  \item 外延层堆栈决定的纵向约束和损耗重叠；
  \item 有限孔阵、边缘吸收和窗口定义决定的横向模式边界。
\end{itemize}

如果只把“周期孔阵”画出来，而不把真实外延层、接触层和边缘终止方式带进去，那么算出的模式常常只是“美观的理想光学对象”，而不是器件会真的支持的模式。也就是说，几何是否足够真实，直接决定你在算“物理问题”还是在算“几何练习题”。

这一点在 PCSEL 上尤其严厉，因为其外延层既参与光学，也参与后续电学和热学接口。光场哪怕只是在纵向上偏移一点，量子阱重叠和掺杂吸收就都可能随之改变。

\section{边界条件：PML 不是附件，而是模型的一部分}
只要求解开放器件，边界条件就是结果的一部分，不是外加包装。对 FDTD 和 FEM 来说，最常见的开放边界处理都是 PML。它的目标是在有限计算域内尽量模拟“波向无穷远走掉而不反射回来”。

但初学者很容易把 PML 当成一个固定配件，仿佛加上去就自然没问题。实际上，PML 厚度、距离器件的间隔、坐标拉伸参数和离散质量，都会直接影响共振频率、$Q$ 值和远场。若 PML 太近，器件近场和倏逝场会直接碰到吸收层；若 PML 太薄或参数不当，又会产生伪反射，污染高 $Q$ 模和远场主瓣。

因此，对任何声称给出辐射 $Q$ 或远场的 FDTD/FEM 结果，都应该追问：边界离器件多远？PML 厚度怎么选？结果对这些参数是否收敛？如果这些问题答不出来，那么结果的可信度就不够。

\begin{table}[htbp]
  \centering
  \footnotesize
  \caption{FDTD/FEM 中开放边界设置的最小量化检查表。数值不是普适常数，而是第一次建模时可执行的起点。}
  \label{tab:pml_quantitative_rules}
  \begin{tabularx}{\textwidth}{@{}p{3.0cm}p{4.0cm}Y@{}}
    \toprule
    项目 & 建议起点 & 通过门槛 \\ \midrule
    PML 与器件最近距离 & 至少取 $d_{\mathrm{PML}}\gtrsim \lambda_0/n_{\mathrm{eff}}$；若存在强倏逝尾则进一步外推 & 共振频率、$Q$ 和远场主瓣对该距离变化小于预设容差 \\
    PML 厚度 & 至少 $L_{\mathrm{PML}}\gtrsim \lambda_0/(2n_{\min})$，且不少于 $8$--$12$ 个网格层 & 继续增厚时，反射谱、$Q$ 与远场变化应进入平台区 \\
    PML 吸收强度 / 拉伸参数 & 先用默认稳定值，再单因子扫描使理论边界反射压到 $10^{-6}$ 量级以下 & 不应出现随吸收参数变化而大幅漂移的伪模或伪反射 \\
    远场监视面到器件距离 & 置于近场主结构之外，避免切到强局域源区 & 改变监视面位置时，角谱主瓣与旁瓣结论不翻转 \\
    \bottomrule
  \end{tabularx}
\end{table}

更具体地说，若要报告一个可追溯的开放边界结果，至少应给出三条收敛句：第一，PML 与器件距离从 $d_{\mathrm{PML}}^{(1)}$ 增加到 $d_{\mathrm{PML}}^{(2)}$ 时，目标模频率与 $Q$ 的变化量；第二，PML 厚度从 $L_{\mathrm{PML}}^{(1)}$ 增加到 $L_{\mathrm{PML}}^{(2)}$ 时的变化量；第三，远场主瓣角度与主瓣半高全宽是否对监视面位置稳定。若这三条信息没有出现，那么“边界已经足够远”通常只是主观判断。

若使用 FDTD，还必须同时给出空间网格与时间步长的起点。对三维 Yee 网格，稳定性的最小 Courant 条件可写成
\begin{equation}
\Delta t
\le
\frac{1}{c}
\left(
\frac{1}{\Delta x^2}
+
\frac{1}{\Delta y^2}
+
\frac{1}{\Delta z^2}
\right)^{-1/2}.
\label{eq:courant_condition_fdtd}
\end{equation}
工程上通常再乘一个安全因子 $S<1$，例如取 $\Delta t=S\Delta t_{\max}$、$S\sim0.5$。同时，空间网格也应至少满足
\[
\Delta_{\min}\lesssim \min\!\left(\frac{\lambda_0}{10n_{\max}},\frac{r_{\mathrm{feature}}}{5}\right),
\]
其中 $r_{\mathrm{feature}}$ 代表最小孔边圆角、最窄脊宽或最薄有效层厚度。对量子阱、金属邻近层和陡侧壁附近，还应再做局部加密。这样写的意义不是要求读者死记一个数字，而是提醒：\textbf{时间步长不是独立自由度，它被最细网格直接锁定；而最细网格又必须由波长分辨率和最小几何特征共同决定。} 只要孔边、薄量子阱附近或金属邻近层要求更细的 $\Delta x,\Delta y,\Delta z$，$\Delta t$ 也必须同步减小。否则，即便频谱看上去平滑，结果也可能只是数值色散和时间离散误差共同伪装出来的“稳定解”。

\section{激励源的选择，其实是在决定你想看什么}
FDTD 和 FEM 并不自动知道你关心哪一支模式。源的放置方式、方向、时域或频域特性，实际上是在告诉求解器“请优先激发哪些对称性和哪些频段”。

如果目标是搜寻候选共振频率，宽带脉冲源通常很合适；如果目标是选择性激发某种偏振或某类对称性模式，则应使用与目标模式近场结构匹配的电流源、偶极源或端口源；如果目标是得到稳态散射响应，则频域连续波源更合适。

这一步之所以重要，是因为错误源设置并不会总是导致仿真报错，而更常见的是：它悄悄让你只看到某几支容易被激发的模式，误以为系统里就只有这些模式。对接近简并的 PCSEL 候选模来说，这种误导尤其危险。

\begin{ExampleBox}
若目标是比较一对近简并正交偏振 Gamma 点模，而激励源只沿单一方向放置且恰好与其中一支模式强耦合，那么仿真中另一支模式可能几乎不出现。此时若不更换源方向或改用本征模分析，就可能错误得出“器件天然单偏振”的结论。
\end{ExampleBox}

\section{怎样从时域或频域结果里提取共振频率与品质因子}
对冷腔分析来说，最常提取的两个量是共振频率和品质因子。FDTD 中，常见做法是在脉冲激励后观察某个监视点的场随时间衰减，再对包络做指数拟合，从而得到衰减率 $\gamma$，进而写成
\begin{equation}
Q=\frac{\omega_0}{2\gamma}.
\label{eq:q_from_decay}
\end{equation}

FEM 中，则常通过开放本征值求解直接得到复频率
\begin{equation}
\tilde\omega=\omega_0-\ii\gamma,
\label{eq:complex_eig_fullwave}
\end{equation}
再由 \cref{eq:q_from_decay} 得到 $Q$。两条路在理想条件下是一致的，但它们对数值误差的敏感方式不同。因此，若条件允许，让时域和频域对同一几何给出一致的趋势，是很有价值的互证。

这里再次要强调：冷腔频率和冷腔 $Q$ 是重要输入，但仍不是器件工作点本身。把它们直接当作激射频率和工作阈值，是层次上的错误。

\section{近场、远场与模场：怎样读出真正有用的量}
很多全波仿真做完之后，会得到大量漂亮图像。真正重要的是知道该从中提取什么。

首先是模场分布。它不是为了“看起来像不像文献图”，而是为了和后续的量子阱重叠、损耗层重叠、偏振结构和边缘泄漏联系起来。

其次是远场。PCSEL 的一个核心性能就是远场与偏振可工程化，因此必须确认远场主瓣、旁瓣和角宽是否对网格、边界和监视面位置稳定。

再次是功率流和损耗分解。如果可能，应分辨功率是主要向上辐射、向下辐射，还是从侧边泄漏。否则即便 $Q$ 给出来了，也难以理解损耗的真正来源。

最后是与器件层的接口量，例如与量子阱区的重叠因子、与重掺杂层或金属层的场重叠。这些量往往比单一“最大场强”更重要，因为它们直接进入 \cref{ch:epitaxial-optics,ch:qw-gain,ch:electrical,ch:eto-coupling} 的器件模型。

若要把近场变成真正可比较的远场，至少要记住角谱或近远场变换的最小形式。对监视面上的切向场，远场复振幅可形式化写成
\begin{equation}
\bm{F}(k_x,k_y)
\propto
\iint_{S_{\mathrm{mon}}}
\bm{E}_t(x,y)\,
\ee^{-\ii(k_x x+k_y y)}\,
\dd x\,\dd y,
\label{eq:angular_spectrum_nf2ff}
\end{equation}
其观察角满足
\begin{equation}
k_x = n_{\mathrm{out}}k_0\sin\theta\cos\phi,
\qquad
k_y = n_{\mathrm{out}}k_0\sin\theta\sin\phi.
\label{eq:angular_coordinates_nf2ff}
\end{equation}
\cref{eq:angular_spectrum_nf2ff,eq:angular_coordinates_nf2ff} 的意义不是要求读者手写近远场变换程序，而是提醒两个常见事实：第一，远场角谱本质上是监视面切向场的二维傅里叶变换；第二，若监视面过近、截面过小或含有强反应近场，角谱会被直接污染。

\begin{PitfallBox}
“谱污染”在这里通常不是指图像不够平滑，而是指你以为看到的是器件的远场或本征谱，实际上看到的是边界反射、监视面截断、时间窗不够长或错误激励源带来的伪结构。只要改变监视面位置、时间窗或 PML 设置后主结论就翻转，就应先怀疑谱污染，而不是先怀疑新物理。
\end{PitfallBox}

\section{冷腔分析和带增益分析到底差在哪里}
对初学者来说，一个特别容易混淆的问题是：在全波方法里“加一点损耗或增益”看起来只是改材料参数，为什么物理上层级却变了？原因是冷腔分析求解的是被动线性系统，而加入增益后，系统开始逼近一个有源开放系统。

如果只是加入一个均匀、小信号的虚部增益并保持线性，那么你得到的仍然只是“光学阈值趋势”或“复频率如何变化”的信息；一旦再加入载流子分布、增益饱和和温升反馈，问题就不再是单纯电磁本征问题，而变成光-电-热耦合问题。也就是说，数学上的变化不仅是“参数值改了”，而是“方程组和闭合条件都改了”。

\begin{PitfallBox}
很多工作中会把“冷腔全波仿真 + 均匀增益补偿到零衰减”的结果直接称为“器件阈值仿真”。这种说法是不严谨的。它最多是光学线性阈值估计，不是电注入器件级阈值电流预测。
\end{PitfallBox}

\begin{table}[htbp]
  \centering
  \footnotesize
  \caption{冷腔全波分析与带增益线性分析的结论边界。}
  \label{tab:cold_vs_gain_fullwave}
  \begin{tabularx}{\textwidth}{@{}p{2.5cm}p{4.5cm}Y@{}}
    \toprule
    层级 & 主要输入/输出 & 不能越级宣称什么 \\ \midrule
    冷腔全波 & 输入为实介电常数、开放边界与有限器件几何；输出为复频率、$Q$、近场、远场、损耗分解 & 不能直接宣称阈值增益、阈值电流或工作频率 \\
    小信号带增益全波 & 在冷腔基础上加入均匀或给定分布的小信号增益代理；输出为复频率趋向、线性阈值趋势、竞争模排序变化 & 不能直接宣称阈上稳态、增益钳位、热稳定工作窗口 \\
    电热光耦合工作点 & 输入为注入、热边界、增益压缩与回写模型；输出为阈值电流、工作频率漂移、模式稳定性 & 仍需实验回验才能宣称器件最终性能 \\
    \bottomrule
  \end{tabularx}
\end{table}

\section{收敛检查：哪些量必须扫，哪些量必须复核}
全波结果是否可信，不能靠肉眼判断。最基本的检查至少包括四类。

第一类是空间离散收敛。FDTD 需要扫描网格尺寸；FEM 需要扫描网格密度和局部加密策略。

第二类是边界收敛。PML 厚度、边界距离和吸收参数必须被扫过，而不能只选一组看上去顺眼的数。

第三类是时间或求解容差收敛。FDTD 中要检查时间步长和总仿真时长；FEM 中要检查求解残差、特征值搜索区间和迭代容差。

第四类是结果提取方式收敛。不同监视点、不同拟合窗口、不同近远场变换面，不应导致结论大幅波动。

如果这四类检查都过不去，那么即使图形非常漂亮，也不能进入器件级推断。

对本科生或第一次上手的研究者，更可执行的说法是：至少固定一份最小收敛账本。FDTD 至少记录 $\Delta x$、Courant 因子、总时间窗 $T_{\mathrm{sim}}$、PML 厚度与距离；FEM 至少记录网格最小尺寸、单元阶次、特征值搜索区间、残差容差和 PML 设定。只有账本存在，后续“结果为什么可信”才不是回忆题。

\section{本章在整条器件主线中的位置}
把全书逻辑接起来，本章起的是承上启下的作用。\cref{ch:rcwa} 给出开放周期单胞层面的结论，而本章把问题推进到有限尺寸全波层面。接下来 \cref{ch:method-compare} 会总结为什么单一方法不能覆盖所有问题，并给出推荐的互证工作流。再往后 \cref{ch:epitaxial-optics,ch:qw-gain,ch:electrical,ch:eto-coupling} 则把这些光学结果转交给外延、增益、电注入和热耦合模型，进入真正的器件主线。

因此，FDTD 和 FEM 在书中的角色，不是简单地“又多了两个求解器”，而是把“候选模光学”转化为“真实器件光学”的关键桥梁。

\section*{本章小结}
FDTD 与 FEM 都是在有限尺寸、开放边界下求解 Maxwell 方程的全波方法，但它们擅长的对象不同。FDTD 长于宽带和瞬态，FEM 长于复杂几何和高精度本征模分析。对 PCSEL 而言，它们的真正价值在于把无限周期单胞结论推进到真实有限器件，并输出与外延、增益、电注入和热耦合模型兼容的场分布与损耗信息。它们都很强，但都不能跳过器件级多物理场闭环而直接宣布最终工作结论。

\section*{练习题}
\begin{enumerate}
  \item \basicq 解释为什么在 RCWA 之后，有限尺寸器件分析不能被简单省略。

  \item \basicq 比较 FDTD 与 FEM 在“宽带扫频”“复杂孔形”“高 $Q$ 本征模”和“瞬态建立过程”四个问题上的优劣。

  \item \midq 设计一套用于 PCSEL 冷腔 $Q$ 提取的收敛检查方案，至少包含网格、边界和结果提取三类参数。

  \item \hardq 给出一个“冷腔全波结果很好看，但器件最终仍会失败”的可能场景，并说明缺失了哪一层物理。
\end{enumerate}

\section*{建议继续阅读}
\begin{itemize}
  \item 下一章将把 CWT、RCWA、FDTD 和 FEM 统一放到同一个工作流里比较。读完下一章后，再回看本章会更清楚每种方法到底负责哪一段问题。 这条阅读的重点是把本章的输出顺着主线接到后续模块，而不是停成孤立知识点。

  \item 若想补充数值方法基础，可参考 \cite{taflove2005,jin2015} 中关于时域方法、有限元方法与开放边界处理的章节，并结合 \cite{berenger1994} 回看 PML 的来历。 这条阅读的重点是补足本章关于 FDTD/FEM、开放边界与有限阵列验证的标准背景、经典语言和代表性文献接口。
\end{itemize}


